% Scheletro di lavoro, non testo definitivo.
% Fonti e passi sono proposte di lettura: verificarli sul testo greco prima della stesura.
% Seguira il metodo: documento originale, traduzione, esame dell'argomentazione, ricostruzione critica.
\chapter{Platone, Socrate e i sofisti}
\label{cap:platone:sofisti}

% Problema: Come distinguere il Socrate storico dal personaggio platonico?
% Fonti da verificare: Apologia 17a–42a; Critone 43a–54e; Eutifrone.
\section{Il problema dell'eredità socratica}

\subsection{Socrate e la ricerca della definizione}

\subsection{Sapere, virtù ed esame della vita}

% Problema: Che rapporto sussiste fra competenza, educazione e virtù?
% Fonti da verificare: Protagora 319a–328d, 349d–361c; Menone 70a–100b.
\section{La virtù si può insegnare?}

\subsection{La tesi di Protagora e le obiezioni di Socrate}

\subsection{Virtù, conoscenza e opinione retta}

% Problema: In che cosa la persuasione può distinguersi da un sapere?
% Fonti da verificare: Gorgia 447a–466a, 481b–527e; Fedro 259e–274b.
\section{Retorica, persuasione e verità}

\subsection{Gorgia, Polo e Callicle}

\subsection{La critica della retorica e la possibilità di un'arte del discorso}

% Problema: Quando un'argomentazione conduce alla verità e quando produce soltanto vittoria dialettica?
% Fonti da verificare: Eutidemo 271a–304c.
\section{L'eristica e i limiti della confutazione}

\subsection{Confutazione socratica ed eristica}

\subsection{Conseguenze educative della differenza}
